Recent work has given particular attention to worst-case analysis. In~\cite{VRBS25}, the authors investigated finite-step convergence verification of~\eqref{eq:first_order_method} for parametric convex quadratic programs (PQPs), obtaining tight, data-free quantifications of $r^k_{\max}(x)$ that incorporate warm-starting effects. Similarly,~\cite{VRBS+24} formulated an exact mixed-integer linear program to maximise $r^k(x)$, defined using the $\infty$-norm, by encoding the algorithm’s dynamics as polyhedral constraints. While these approaches yield rigorous guarantees, they tend to be conservative for typical parameter instances, and they scale poorly when extended to large-scale or non-convex problems. Although less pessimistic than typical computer-assisted methods, purely worst-case bounds can still underestimate algorithmic reliability in real-world applications.  

To address these limitations, probabilistic analyses have emerged as a promising alternative. For instance,~\cite{RSBS25} developed probabilistic guarantees for both classical and learned optimisers using sample convergence bounds and the Probably Approximately Correct (PAC)-Bayes framework, respectively. Inspired by this work, the authors of~\cite{JHKM+25} employed advancements in the so-called Scenario Approach Theory to generate tight, data-driven guarantees for both convergence and proximity to the optimal solution through a general, user-defined metric. These works illustrate the potential of probabilistic methods to offer less conservative and more informative guarantees than worst-case analysis. However, existing approaches... \{CONCLUDE BY HIGHLITING DISADVANTAGES e.g. unintuitiveness of PAC Bayes or Sample count-dependence of scenario, I don't know if these are strong enough\}